\documentclass[11pt]{article}
\usepackage[a4paper,margin=25mm]{geometry}
\usepackage{amsmath,amssymb,amsthm,booktabs}
\usepackage[british]{babel}
\usepackage{url,xcolor,graphicx}
\usepackage[round,authoryear]{natbib}
\usepackage[hidelinks]{hyperref}
\setlength{\bibhang}{1.5em}
\setlength{\bibsep}{2pt}
\setlength{\emergencystretch}{1.5em}
\renewcommand{\bibfont}{\small}
\renewcommand{\topfraction}{0.85}
\renewcommand{\textfraction}{0.1}
\renewcommand{\floatpagefraction}{0.8}
\newtheorem{theorem}{Theorem}[section]
\newtheorem{proposition}[theorem]{Proposition}
\title{Capped GBM pricing\\[2mm]\large Replication, numerical convergence and Monte Carlo sensitivities}
\author{Amanjeet Singh}
\date{11 October 2026}
\begin{document}
\maketitle
\begin{abstract}
I connect no-arbitrage replication, martingale valuation and numerical
pricing for a capped European claim under geometric Brownian motion with
continuous dividends. A heat-coordinate transformation permits a monotone
backward-Euler scheme with unconditional maximum-norm stability. I prove
convergence for the actual kinked payoff and separate finite-domain,
interpolation and discretisation errors. Explicit central-difference bounds
establish joint price and Greek consistency when domains, meshes and bumps
are refined appropriately. Exact terminal simulation supplies pathwise
first sensitivities and likelihood-ratio gamma. A mixed sensitivity
derivation gives a bounded gamma estimator with finite-sample concentration.
Seven synthetic contracts compare analytic, PDE and Monte Carlo prices and
five Greeks at declared sampling budgets. The mixed gamma estimator lowers
measured raw variance, while independently fitted stock controls worsen
some sensitivities. Detailed derivations and adverse numerical findings are
retained. The analytic call spread remains the strongest valuation baseline;
this companion demonstrates and checks the financial and computational
foundations supporting the basket-pricing project.
\end{abstract}
\noindent\begin{minipage}{\linewidth}
\small\emergencystretch=2em
\textit{Keywords:} derivative pricing; geometric Brownian motion; replication;
finite differences; numerical stability; convergence; Monte Carlo Greeks.\par
\smallskip\textit{JEL classification:} C15, C63, G13.\par
\end{minipage}
% The parent document supplies amsmath, amssymb, amsthm, natbib and hyperref.
% This fragment deliberately contains no claim of theoretical originality.

\section{Financial model and no-arbitrage valuation}
\label{sec:cgbm-model}

I use a one-asset claim to connect replication, martingale valuation,
parabolic equations and statistical estimation under one financial model.
The purpose is to demonstrate the mathematics underlying the numerical
comparisons. The classical valuation framework follows
\citet{blackscholes1973} and \citet{merton1973}; these established methods
do not establish a quantum advantage or a new financial theory. Replication and call valuation specialise
\citet{blackscholes1973}; the continuous-dividend treatment follows
\citet{merton1973}. The pricing and sensitivity derivations below retain
the assumptions required for these specialisations.

Let $S>0$, $K>0$, $H>0$, $\sigma>0$ and $\tau>0$. The interest rate $r$ and
continuous dividend yield $q$ are constant. Write $J=K+H$ and
\begin{equation}
 F(x)=\min\{(x-K)^+,H\}=(x-K)^+-(x-J)^+,
 \qquad D=\exp(-r\tau),\qquad v=\sigma\sqrt{\tau}.
 \label{eq:cgbm-payoff}
\end{equation}
Under the selected pricing measure $Q$, the ex-dividend asset follows
\begin{equation}
 dS_t=(r-q)S_t\,dt+\sigma S_t\,dW_t^Q,
 \qquad
 X=S\exp\{b\tau+vZ\},
 \qquad b=r-q-\tfrac12\sigma^2,
 \label{eq:cgbm-law}
\end{equation}
where $Z$ is standard normal. The price is
\begin{equation}
 V(S,\tau;r,q,\sigma)=D\,\mathbb E_Q[F(X)].
 \label{eq:cgbm-price}
\end{equation}
This is a model comparison with synthetic parameters, rather than a
calibrated empirical test. A physical return mean cannot be substituted
for the pricing drift $r-q$.
I assume frictionless continuous trading, reinvestable dividends,
borrowing and lending at the same rate $r$, and an unrestricted stock
position within the stated admissible strategy class. The probability
space carries the usual augmented Brownian filtration, satisfying the
usual completeness and right-continuity conditions.

\begin{proposition}[Replication and the pricing equation]
\label{prop:cgbm-replication}
The price in \eqref{eq:cgbm-price} is bounded, continuous at expiry and smooth
for $S,\tau>0$. It satisfies
\begin{equation}
 V_\tau=\tfrac12\sigma^2S^2V_{SS}+(r-q)SV_S-rV,
 \qquad V(S,0)=F(S).
 \label{eq:cgbm-pde}
\end{equation}
Holding $V_S$ shares and placing $V-SV_S$ in the money-market account
replicates the claim when dividends are reinvested within the
self-financing strategy. Its discounted wealth is non-negative, and its
stochastic gain is square-integrable on each finite horizon. Under the
usual admissibility condition that discounted wealth is bounded below,
the stipulated market has no arbitrage. This replicating price is the
unique no-arbitrage price among bounded claim valuations and the unique
bounded classical solution of the terminal-value problem. Here a classical
solution is $C^{1,2}$ for positive remaining maturity and is continuous
at expiry.
\end{proposition}

The dividend term belongs in the self-financing gain equation, not in
the terminal payoff. Appendix~\ref{app:cgbm-financial} derives both the
measure change and the replication argument, including the limiting
argument at the payoff kinks.

\section{Analytic prices and Greeks}
\label{sec:cgbm-analytic}

Let $\Phi$ and $\varphi$ denote the standard normal distribution function
and density. For strike $A>0$, define
\begin{equation}
 d_{1,A}=\frac{\log(S/A)+(r-q+\sigma^2/2)\tau}{v},
 \qquad d_{2,A}=d_{1,A}-v,
 \qquad C_A=S e^{-q\tau}\Phi(d_{1,A})-A D\Phi(d_{2,A}).
 \label{eq:cgbm-call}
\end{equation}
Then $V=C_K-C_J$. Write
$\Delta\Phi_i=\Phi(d_{i,K})-\Phi(d_{i,J})$ and
$\Delta\varphi=\varphi(d_{1,K})-\varphi(d_{1,J})$. Differentiation gives
\begin{align}
 \Delta=V_S&=e^{-q\tau}\Delta\Phi_1,\label{eq:cgbm-delta}\\
 \Gamma=V_{SS}&=\frac{e^{-q\tau}\Delta\varphi}{Sv},\label{eq:cgbm-delta-gamma}\\
 \mathcal V=V_\sigma&=S e^{-q\tau}\sqrt{\tau}\,\Delta\varphi,\label{eq:cgbm-vega}\\
 \mathcal R=V_r&=\tau D\{K\Phi(d_{2,K})-J\Phi(d_{2,J})\},\label{eq:cgbm-vega-rho}\\
 \Theta=-V_\tau&=-\frac{S e^{-q\tau}\sigma}{2\sqrt{\tau}}\Delta\varphi
 +q S e^{-q\tau}\Delta\Phi_1\notag\\
 &\quad-rD\{K\Phi(d_{2,K})-J\Phi(d_{2,J})\}.\label{eq:cgbm-theta}
\end{align}
Vega and rho are derivatives per unit change in volatility and interest
rate. Calendar theta uses a fixed expiry and advances the valuation date,
so it is the negative maturity derivative. Gamma and vega can be negative
for a capped claim; this is a property of the call spread, not a numerical
failure.

\section{A monotone heat-equation discretisation}
\label{sec:cgbm-heat}

Introduce $y=\log(S/K)+b\tau$, $a=\sigma^2/2$ and
$f(y)=F(K e^y)$. The whole-line heat solution $w$ satisfies
\begin{equation}
 w_\tau=a w_{yy},\qquad w(0,y)=f(y),\qquad
 V(S,\tau)=D w(\tau,y).
 \label{eq:cgbm-heat}
\end{equation}
Its payoff kinks lie at $0$ and $c=\log(J/K)$. Moreover,
$0\leq f\leq H$ and $f$ is globally Lipschitz with constant
$L_f=J$. For a finite domain $L<0<c<U$, let $w^{\mathrm{box}}$ solve the
same equation with fixed boundary values $0$ and $H$.
These values approximate the exact financial far-field limits
\begin{equation}
 \lim_{S\downarrow0}V(S,\tau)=0,\qquad
 \lim_{S\to\infty}V(S,\tau)=DH.
 \label{eq:cgbm-far-field}
\end{equation}
Indeed the terminal asset tends pointwise to zero or infinity as its
initial spot does, and the bounded continuous payoff permits dominated
convergence. At remaining maturity $\tau$, the heat-domain endpoints
correspond to asset values $K\exp(L-b\tau)$ and
$K\exp(U-b\tau)$. Their finite values do not make the imposed boundary
prices exact; the exit-tail bound below quantifies the approximation.

For integers $M\geq3$ and $N\geq1$, set
$h=(U-L)/M$, $y_j=L+jh$, $k=\tau/N$ and
$\lambda=ak/h^2$. Backward Euler gives
\begin{equation}
 -\lambda W_{j-1}^{n}+(1+2\lambda)W_j^{n}
 -\lambda W_{j+1}^{n}=W_j^{n-1},\qquad
 W_0^n=0,\quad W_M^n=H,\quad W_j^0=f(y_j).
 \label{eq:cgbm-be}
\end{equation}
I interpolate the final grid values linearly at the query $y$ and then
apply the discount factor. Subtracting the affine boundary lift diagonalises
this same finite-difference operator in the discrete sine basis. Its
eigenvalues and backward-Euler factors are
\begin{equation}
 \nu_\ell=\frac{4}{h^2}\sin^2\!\left(\frac{\ell\pi}{2M}\right),
 \qquad (1+ak\nu_\ell)^{-N},\qquad 1\leq\ell\leq M-1.
 \label{eq:cgbm-dst}
\end{equation}
This is a faster evaluation of \eqref{eq:cgbm-be}, rather than a different
spatial method. A direct tridiagonal time-step implementation checks that
equivalence numerically. A floating-point implementation is not an
interval-arithmetic certificate merely because its mathematical scheme
has a proved bound.

\begin{theorem}[Unconditional stability and a nonsmooth-payoff error bound]
\label{thm:cgbm-be}
For every $h,k>0$, \eqref{eq:cgbm-be} preserves $0\leq W_j^n\leq H$.
Differences between two solutions with the same boundary values contract
in the maximum norm at every step. Fix an even, non-negative
$\rho\in C_c^\infty(-1,1)$ with integral one, and set
\begin{equation}
 m_1=\int_{-1}^1|z|\rho(z)\,dz,\qquad
 c_3=\|\rho'''\|_{L^1},\qquad
 0<\eta<\min\{-L,U-c\}.
 \label{eq:cgbm-mollifier}
\end{equation}
For exact arithmetic, the nodal error at maturity obeys
\begin{equation}
 \max_{1\leq j\leq M-1}
 |W_j^N-w^{\mathrm{box}}(\tau,y_j)|
 \leq 2L_fm_1\eta
 +\tau\left(\frac{ah^2}{12}+\frac{a^2k}{2}\right)
 L_f c_3\eta^{-3}.
 \label{eq:cgbm-be-error}
\end{equation}
Consequently the scheme converges to the finite-domain solution, with
a rigorous upper bound of order $(h^2+k)^{1/4}$ for this nonsmooth payoff
on a fixed time interval. This conservative bound does not assert that
the observed rate is exactly one quarter.
\end{theorem}

The proof uses the maximum principle, a compact mollifier and an odd
periodic extension after the boundary lift. It does not assume that the
unmollified payoff has four derivatives. The full proof is in
Appendix~\ref{app:cgbm-numerical}.

\begin{proposition}[Domain and interpolation errors]
\label{prop:cgbm-domain}
For $x\in[L,U]$, put
\begin{equation}
 \beta(\tau,x)=\min\!\left\{1,
 2\Phi\!\left(-\frac{x-L}{\sigma\sqrt{\tau}}\right)
 +2\Phi\!\left(-\frac{U-x}{\sigma\sqrt{\tau}}\right)\right\}.
 \label{eq:cgbm-domain-beta}
\end{equation}
Then $|w(\tau,x)-w^{\mathrm{box}}(\tau,x)|\leq H\beta(\tau,x)$.
Let $y$ lie between grid nodes $y_j,y_{j+1}$ and define
$\beta_* =\max\{\beta(\tau,y_j),\beta(\tau,y_{j+1})\}$.
The interpolated, discounted price $\widehat V$ satisfies
\begin{align}
 |\widehat V-V|\leq D\Bigg[
 &H\beta_*+2L_fm_1\eta
 +\tau\left(\frac{ah^2}{12}+\frac{a^2k}{2}\right)
 L_f c_3\eta^{-3}\notag\\
 &+\frac{h^2 L_f}{8\sigma\sqrt{\tau}}
 \sqrt{\frac{2}{\pi}}\Bigg].
 \label{eq:cgbm-total-error}
\end{align}
\end{proposition}

The four terms control the finite domain, payoff smoothing, numerical
consistency and interpolation. Payoff smoothing is a proof device; the
implemented initial condition remains the original capped payoff.
On a fixed domain, mesh refinement converges to
$w^{\mathrm{box}}$, leaving a domain error relative to the whole-line
financial price. Full-price consistency requires the domain error to
vanish as well.

\section{Greek consistency and Monte Carlo comparisons}
\label{sec:cgbm-greeks}

Nonsmooth payoffs require separate control of sensitivities; the analysis of 
\citet{gilescarter2006} gives a published comparison for Crank--Nicolson/
Rannacher schemes. I prove the bounds for the backward-Euler scheme used here.

I compute PDE Greeks by central price differences on the same domain and
grid. This fixes the numerical heat domain under a bump; the transformed
query $y$, diffusivity $a$ and discount factor still change with the bumped
parameter. The small-bump limit must respect the price error.

\begin{proposition}[Propagation into central differences]
\label{prop:cgbm-bumps}
Suppose $|\widehat V-V|\leq E$ at every required bumped point. For a scalar
parameter $p$, if $|\partial_p^3 V|\leq M_3$ on its bump interval, then
\begin{equation}
 \left|\frac{\widehat V(p+d)-\widehat V(p-d)}{2d}
 -\partial_p V(p)\right|
 \leq \frac{E}{d}+\frac{M_3d^2}{6}.
 \label{eq:cgbm-first-bump}
\end{equation}
For spot gamma, if $|\partial_S^4 V|\leq M_4$ on the spot bump interval,
then
\begin{equation}
 \left|\frac{\widehat V(S+d)-2\widehat V(S)+\widehat V(S-d)}{d^2}
 -\Gamma\right|
 \leq \frac{4E}{d^2}+\frac{M_4d^2}{12}.
 \label{eq:cgbm-second-bump}
\end{equation}
These statements apply uniformly on a compact parameter set with
$S,\sigma,\tau$ bounded away from zero, provided all queries lie inside
the domain. Refining the mesh and enlarging the domain so that $E_n\to0$,
then choosing $d_n\to0$ and $E_n/d_n\to0$ for first derivatives or
$E_n/d_n^2\to0$ for gamma, proves consistency for the whole-line Greeks.
\end{proposition}

After fixing units, choices proportional to $E_n^{1/3}$ and $E_n^{1/4}$
balance the displayed error terms for first derivatives and gamma,
respectively. Practical comparisons also retain the exact analytic
central-difference bias at each implemented bump. This diagnoses numerical
price error separately from bump bias and prevents an apparent Greek
discrepancy from being assigned wholly to the PDE scheme.

The implemented bump recipe uses a dimensionless log-grid spacing $h$,
a positive spot reference scale $S_*$, and positive reference scales for
the other parameters. The spot bump is proportional to $S_*h^{1/8}$;
first-order parameter bumps are proportional to $h^{1/4}$ in their
respective reference units. With $k=O(h^2)$, the rigorous price bound is
$O(h^{1/2})$ after making the domain error $o(h^{1/2})$. The resulting
spot gamma error is $O(h^{1/4})$, and the first-order errors are at most
$O(h^{1/4})$. These are conservative guaranteed upper rates, not fitted
empirical orders. A finite fixed domain retains its displayed exit-tail
allowance, regardless of this bump recipe.

Pathwise and likelihood-ratio estimators have established foundations
\citep{broadieglasserman1996}. The capped-payoff formulae and their ranges
are derived explicitly in Appendix~\ref{app:cgbm-monte-carlo}.

For exact terminal sampling, write $I=\mathbf 1_{\{K<X<J\}}$ and
define the following one-path estimators:
\begin{align}
 Y_V&=D F(X),
 &Y_\Delta&=D I X/S,\notag\\
 Y_{\mathcal V}&=D I X(\sqrt{\tau}Z-\sigma\tau),
 &Y_{\mathcal R}&=D\tau\{IX-F(X)\},
 \label{eq:cgbm-pw-estimators}\\
 Y_\Theta&=D\left[rF(X)-IX
 \left(b+\frac{\sigma Z}{2\sqrt{\tau}}\right)\right],
 \label{eq:cgbm-theta-estimator}\\
 Y_\Gamma&=\frac{D F(X)}{S^2}
 \left\{\frac{Z^2-1}{v^2}-\frac{Z}{v}\right\}.
 \label{eq:cgbm-lr-estimator}
\end{align}
The first-order Greeks use pathwise differentiation, following the
simulation sensitivity framework of \citet{broadieglasserman1996}.
Gamma uses the
second likelihood-ratio weight: differentiating the pathwise indicator
as though it were smooth would omit its boundary contribution.
I also retain the bounded mixed pathwise/likelihood-ratio rival
\begin{equation}
 Y_{\Gamma,\mathrm{mix}}
 =\frac{D I X}{S^2}\left(\frac{Z}{v}-1\right).
 \label{eq:cgbm-mixed-gamma}
\end{equation}
It follows by differentiating delta's fixed-asset density representation,
rather than differentiating a random indicator pointwise. Its derivation
appears in Appendix~\ref{app:cgbm-monte-carlo}. Boundedness gives a
different finite-sample uncertainty guarantee, but does not by itself
prove variance dominance over the ordinary likelihood-ratio estimator.

\begin{proposition}[Unbiased estimators and finite-sample uncertainty]
\label{prop:cgbm-mc}
Each estimator in \eqref{eq:cgbm-pw-estimators}--\eqref{eq:cgbm-mixed-gamma}
has expectation equal to its analytic price or Greek and finite second
moment. Let
\begin{equation}
 z_- =\frac{\log(K/S)-b\tau}{v},\quad
 z_+ =\frac{\log(J/S)-b\tau}{v},\quad
 m=\max\{|z_-|,|z_+|\}.
 \label{eq:cgbm-active-z}
\end{equation}
Deterministic ranges for price, delta, vega, rho and theta are, respectively,
\begin{align}
 &[0,DH],\quad[0,DJ/S],\quad[-B_{\mathcal V},B_{\mathcal V}],
 \quad[-D\tau H,D\tau K],\quad[-B_\Theta,B_\Theta],\notag\\
 &B_{\mathcal V}=DJ(\sqrt{\tau}m+\sigma\tau),\qquad
 B_\Theta=D\left[|r|H+J\left(|b|+
 \frac{\sigma m}{2\sqrt{\tau}}\right)\right].
 \label{eq:cgbm-mc-ranges}
\end{align}
For $n$ independent raw paths and a range width $R$, a distribution-free
two-sided confidence radius with failure probability at most $\alpha$,
following \citet{hoeffding1963}, is
\begin{equation}
 R\sqrt{\frac{\log(2/\alpha)}{2n}}.
 \label{eq:cgbm-hoeffding}
\end{equation}
For gamma, the bound
\begin{equation}
 \operatorname{Var}(Y_\Gamma)
 \leq B_\Gamma^2
 :=\frac{D^2H^2}{S^4}\left(\frac{2}{v^4}+\frac{1}{v^2}\right)
 \label{eq:cgbm-gamma-variance}
\end{equation}
gives the finite-sample Chebyshev radius $B_\Gamma/\sqrt{n\alpha}$.
The mixed gamma estimator has range
\begin{equation}
 [-B_{\mathrm{mix}},B_{\mathrm{mix}}],\qquad
 B_{\mathrm{mix}}=\frac{DJ}{S^2}\left(\frac{m}{v}+1\right),
 \label{eq:cgbm-mixed-gamma-bound}
\end{equation}
so it admits the Hoeffding radius with width $2B_{\mathrm{mix}}$.
Allocating failure probabilities across outputs and using a union bound
gives a simultaneous statement. Sample-variance normal intervals are
separate asymptotic diagnostics and do not have these finite-sample
coverage guarantees.
\end{proposition}

The numerical study allocates $0.05/49$ to each of the six ordinary raw
outputs and the additional raw mixed gamma across seven contracts.
The corresponding 49-interval family has failure probability at most
0.05 under the ideal independent-normal sampling model. Comparisons
between Monte Carlo methods additionally report sample-variance normal
intervals as asymptotic diagnostics.

The strongest implemented Monte Carlo comparison combines antithetic
normal pairs with a centred discounted-stock control,
$D X-S e^{-q\tau}$. I fit each control coefficient on independent pilot
pairs and hold it fixed for the main sample. Conditional on the pilot,
the controlled mean is unbiased because the control has expectation zero.
Each independent pair, rather than its two correlated paths, is the
sampling unit. The pilot and both terminal evaluations per pair belong
in the computational work. The controlled estimators contain an unbounded
lognormal control, so their sample-variance intervals remain asymptotic;
the bounded raw-path Hoeffding guarantees are not transferred to them.

\section{Synthetic comparison design and numerical evidence}
\label{sec:cgbm-results}

I use seven designed contracts, all with $K=100$, to exercise both payoff
kinks, the interior spread, deep capping, a short low-volatility horizon,
negative interest with dividends and a volatile long horizon.
Table~\ref{tab:cgbm-contracts} defines the inputs. No observations are
estimated from LSEG or another data provider. The full comparisons and
source/code hashes are retained in \texttt{results/capped\_gbm}.

\begin{table}[tb]\centering\small
\begin{tabular}{lrrrrrr}\toprule
Contract & $S$ & $H$ & $\tau$ & $\sigma$ & $r$ & $q$\\\midrule
Lower kink & 100 & 20 & 1 & 0.2 & 0.03 & 0.01 \\
Upper kink & 120 & 20 & 1 & 0.2 & 0.03 & 0.01 \\
Inside spread & 110 & 20 & 1 & 0.2 & 0.03 & 0.01 \\
Deep cap & 160 & 20 & 1 & 0.2 & 0.03 & 0.01 \\
Low-volatility short dated & 100 & 20 & 0.05 & 0.06 & 0.03 & 0.01 \\
Negative rate/dividend & 105 & 20 & 1 & 0.25 & -0.015 & 0.045 \\
High volatility & 95 & 30 & 2 & 0.65 & 0.03 & 0.01 \\
\bottomrule\end{tabular}
\caption{Wholly synthetic inputs. Rates, volatility and dividend yields are expressed per year; the cap and spot use common monetary units. The labels describe designed test cases rather than an empirical asset sample.}
\label{tab:cgbm-contracts}\end{table}

The coupled meshes use $M=128,256,512,1024,2048,4096,8192,16384$ and
$N=\lceil M^2/8\rceil$. Additional experiments refine time at fixed space
and space at fixed time. The per-contract heat boxes contain both kinks
and use nine diffusion standard deviations of margin, with a minimum
margin of $0.3$. These finite experiments keep the box fixed; the
whole-line convergence theorem additionally requires its tail allowance
to vanish. Domain-enlargement checks preserve both $h$ and $k$, so a
boundary change is not confused with changing the spatial resolution.

The spot bump is
$\min\{0.1S,S(0.006h^{1/8}+h)\}$.
Volatility, rate and maturity bumps are respectively
$\min\{0.25\sigma,0.10\sigma h^{1/4}\}$,
$0.01h^{1/4}$ and
$\min\{0.25\tau,0.05\tau h^{1/4}\}$.
All actual steps and signed bump biases are retained per comparison.
The analytic central difference uses precisely the PDE bump and therefore
isolates its bias. The tridiagonal and sine-basis solvers are compared
at one, three and 64 time steps, including large values of $\lambda$.

\begin{table}[tb]\centering\small
\begin{tabular}{lrl}\toprule
Quantity & Maximum absolute error at $M=16384$ & Units\\\midrule
Price & $9.950\!\times\!10^{-7}$ & Monetary units \\
Delta & $1.111\!\times\!10^{-4}$ & Per unit spot \\
Gamma & $4.137\!\times\!10^{-4}$ & Per unit spot squared \\
Vega & $1.829\!\times\!10^{-3}$ & Per unit volatility \\
Rho & $3.171\!\times\!10^{-4}$ & Per unit rate \\
Theta & $1.026\!\times\!10^{-3}$ & Per year of calendar time \\
\bottomrule\end{tabular}
\caption{Worst absolute discrepancy from the analytic price/Greek across the seven synthetic contracts. Greek errors include the declared bump bias. Different derivative units are not combined into a single accuracy score.}
\label{tab:cgbm-fine-errors}\end{table}

\begin{figure}[tb]
\centering
\setlength{\unitlength}{1mm}
\resizebox{\linewidth}{!}{%
\begin{picture}(164,71)%
\put(14,62){\makebox(61,5)[c]{\small (a) Absolute price error}}%
\put(14.000,22){\line(0,1){38}}%
\put(14.000,17){\makebox(0,3){\scriptsize 128}}%
{\color[HTML]{E5E5E5}\put(22.714,22){\line(0,1){38}}}%
{\color[HTML]{E5E5E5}\put(31.429,22){\line(0,1){38}}}%
\put(31.429,17){\makebox(0,3){\scriptsize 512}}%
{\color[HTML]{E5E5E5}\put(40.143,22){\line(0,1){38}}}%
{\color[HTML]{E5E5E5}\put(48.857,22){\line(0,1){38}}}%
\put(48.857,17){\makebox(0,3){\scriptsize 2048}}%
{\color[HTML]{E5E5E5}\put(57.571,22){\line(0,1){38}}}%
{\color[HTML]{E5E5E5}\put(66.286,22){\line(0,1){38}}}%
\put(66.286,17){\makebox(0,3){\scriptsize 8192}}%
\put(75.000,22){\line(0,1){38}}%
\put(75.000,17){\makebox(0,3){\scriptsize 16384}}%
{\color[HTML]{E5E5E5}\put(14.000,22.000){\line(1,0){61}}}%
\put(12.000,20.500){\makebox(0,3)[r]{\scriptsize $10^{-8}$}}%
{\color[HTML]{E5E5E5}\put(14.000,28.333){\line(1,0){61}}}%
\put(12.000,26.833){\makebox(0,3)[r]{\scriptsize $10^{-7}$}}%
{\color[HTML]{E5E5E5}\put(14.000,34.667){\line(1,0){61}}}%
\put(12.000,33.167){\makebox(0,3)[r]{\scriptsize $10^{-6}$}}%
{\color[HTML]{E5E5E5}\put(14.000,41.000){\line(1,0){61}}}%
\put(12.000,39.500){\makebox(0,3)[r]{\scriptsize $10^{-5}$}}%
{\color[HTML]{E5E5E5}\put(14.000,47.333){\line(1,0){61}}}%
\put(12.000,45.833){\makebox(0,3)[r]{\scriptsize $10^{-4}$}}%
{\color[HTML]{E5E5E5}\put(14.000,53.667){\line(1,0){61}}}%
\put(12.000,52.167){\makebox(0,3)[r]{\scriptsize $10^{-3}$}}%
{\color[HTML]{E5E5E5}\put(14.000,60.000){\line(1,0){61}}}%
\put(12.000,58.500){\makebox(0,3)[r]{\scriptsize $10^{-2}$}}%
\put(14,13){\makebox(61,3)[c]{\scriptsize Space intervals $M$}}%
{\color[HTML]{0072B2}\linethickness{0.5pt}\qbezier(14.000,58.723)(18.357,55.865)(22.714,53.006)}%
{\color[HTML]{0072B2}\linethickness{0.5pt}\qbezier(22.714,53.006)(27.071,49.345)(31.429,45.683)}%
{\color[HTML]{0072B2}\linethickness{0.5pt}\qbezier(31.429,45.683)(35.786,42.675)(40.143,39.667)}%
{\color[HTML]{0072B2}\linethickness{0.5pt}\qbezier(40.143,39.667)(44.500,39.188)(48.857,38.708)}%
{\color[HTML]{0072B2}\linethickness{0.5pt}\qbezier(48.857,38.708)(53.214,35.349)(57.571,31.990)}%
{\color[HTML]{0072B2}\linethickness{0.5pt}\qbezier(57.571,31.990)(61.929,32.270)(66.286,32.551)}%
{\color[HTML]{0072B2}\linethickness{0.5pt}\qbezier(66.286,32.551)(70.643,32.055)(75.000,31.560)}%
{\color[HTML]{0072B2}\put(14.000,58.723){\circle*{0.8}}}%
{\color[HTML]{0072B2}\put(22.714,53.006){\circle*{0.8}}}%
{\color[HTML]{0072B2}\put(31.429,45.683){\circle*{0.8}}}%
{\color[HTML]{0072B2}\put(40.143,39.667){\circle*{0.8}}}%
{\color[HTML]{0072B2}\put(48.857,38.708){\circle*{0.8}}}%
{\color[HTML]{0072B2}\put(57.571,31.990){\circle*{0.8}}}%
{\color[HTML]{0072B2}\put(66.286,32.551){\circle*{0.8}}}%
{\color[HTML]{0072B2}\put(75.000,31.560){\circle*{0.8}}}%
{\color[HTML]{0072B2}\put(14.000,8){\line(1,0){4}}}\put(19.000,6.5){\makebox(0,3)[l]{\scriptsize Lower kink}}%
{\color[HTML]{D55E00}\linethickness{0.5pt}\qbezier(14.000,59.561)(18.357,56.515)(22.714,53.469)}%
{\color[HTML]{D55E00}\linethickness{0.5pt}\qbezier(22.714,53.469)(27.071,51.001)(31.429,48.533)}%
{\color[HTML]{D55E00}\linethickness{0.5pt}\qbezier(31.429,48.533)(35.786,46.170)(40.143,43.808)}%
{\color[HTML]{D55E00}\linethickness{0.5pt}\qbezier(40.143,43.808)(44.500,42.554)(48.857,41.300)}%
{\color[HTML]{D55E00}\linethickness{0.5pt}\qbezier(48.857,41.300)(53.214,37.178)(57.571,33.057)}%
{\color[HTML]{D55E00}\linethickness{0.5pt}\qbezier(57.571,33.057)(61.929,33.850)(66.286,34.644)}%
{\color[HTML]{D55E00}\linethickness{0.5pt}\qbezier(66.286,34.644)(70.643,33.475)(75.000,32.306)}%
{\color[HTML]{D55E00}\put(14.000,59.561){\circle*{0.8}}}%
{\color[HTML]{D55E00}\put(22.714,53.469){\circle*{0.8}}}%
{\color[HTML]{D55E00}\put(31.429,48.533){\circle*{0.8}}}%
{\color[HTML]{D55E00}\put(40.143,43.808){\circle*{0.8}}}%
{\color[HTML]{D55E00}\put(48.857,41.300){\circle*{0.8}}}%
{\color[HTML]{D55E00}\put(57.571,33.057){\circle*{0.8}}}%
{\color[HTML]{D55E00}\put(66.286,34.644){\circle*{0.8}}}%
{\color[HTML]{D55E00}\put(75.000,32.306){\circle*{0.8}}}%
{\color[HTML]{D55E00}\put(46.000,8){\line(1,0){4}}}\put(51.000,6.5){\makebox(0,3)[l]{\scriptsize Upper kink}}%
{\color[HTML]{009E73}\linethickness{0.5pt}\qbezier(14.000,50.759)(18.357,53.165)(22.714,55.572)}%
{\color[HTML]{009E73}\linethickness{0.5pt}\qbezier(22.714,55.572)(27.071,51.773)(31.429,47.974)}%
{\color[HTML]{009E73}\linethickness{0.5pt}\qbezier(31.429,47.974)(35.786,47.221)(40.143,46.469)}%
{\color[HTML]{009E73}\linethickness{0.5pt}\qbezier(40.143,46.469)(44.500,45.572)(48.857,44.675)}%
{\color[HTML]{009E73}\linethickness{0.5pt}\qbezier(48.857,44.675)(53.214,43.009)(57.571,41.342)}%
{\color[HTML]{009E73}\linethickness{0.5pt}\qbezier(57.571,41.342)(61.929,37.264)(66.286,33.186)}%
{\color[HTML]{009E73}\linethickness{0.5pt}\qbezier(66.286,33.186)(70.643,30.244)(75.000,27.302)}%
{\color[HTML]{009E73}\put(14.000,50.759){\circle*{0.8}}}%
{\color[HTML]{009E73}\put(22.714,55.572){\circle*{0.8}}}%
{\color[HTML]{009E73}\put(31.429,47.974){\circle*{0.8}}}%
{\color[HTML]{009E73}\put(40.143,46.469){\circle*{0.8}}}%
{\color[HTML]{009E73}\put(48.857,44.675){\circle*{0.8}}}%
{\color[HTML]{009E73}\put(57.571,41.342){\circle*{0.8}}}%
{\color[HTML]{009E73}\put(66.286,33.186){\circle*{0.8}}}%
{\color[HTML]{009E73}\put(75.000,27.302){\circle*{0.8}}}%
{\color[HTML]{009E73}\put(14.000,4){\line(1,0){4}}}\put(19.000,2.5){\makebox(0,3)[l]{\scriptsize Short dated}}%
{\color[HTML]{CC79A7}\linethickness{0.5pt}\qbezier(14.000,58.452)(18.357,57.875)(22.714,57.298)}%
{\color[HTML]{CC79A7}\linethickness{0.5pt}\qbezier(22.714,57.298)(27.071,54.178)(31.429,51.059)}%
{\color[HTML]{CC79A7}\linethickness{0.5pt}\qbezier(31.429,51.059)(35.786,51.974)(40.143,52.888)}%
{\color[HTML]{CC79A7}\linethickness{0.5pt}\qbezier(40.143,52.888)(44.500,50.223)(48.857,47.557)}%
{\color[HTML]{CC79A7}\linethickness{0.5pt}\qbezier(48.857,47.557)(53.214,42.496)(57.571,37.434)}%
{\color[HTML]{CC79A7}\linethickness{0.5pt}\qbezier(57.571,37.434)(61.929,37.659)(66.286,37.885)}%
{\color[HTML]{CC79A7}\linethickness{0.5pt}\qbezier(66.286,37.885)(70.643,36.269)(75.000,34.653)}%
{\color[HTML]{CC79A7}\put(14.000,58.452){\circle*{0.8}}}%
{\color[HTML]{CC79A7}\put(22.714,57.298){\circle*{0.8}}}%
{\color[HTML]{CC79A7}\put(31.429,51.059){\circle*{0.8}}}%
{\color[HTML]{CC79A7}\put(40.143,52.888){\circle*{0.8}}}%
{\color[HTML]{CC79A7}\put(48.857,47.557){\circle*{0.8}}}%
{\color[HTML]{CC79A7}\put(57.571,37.434){\circle*{0.8}}}%
{\color[HTML]{CC79A7}\put(66.286,37.885){\circle*{0.8}}}%
{\color[HTML]{CC79A7}\put(75.000,34.653){\circle*{0.8}}}%
{\color[HTML]{CC79A7}\put(46.000,4){\line(1,0){4}}}\put(51.000,2.5){\makebox(0,3)[l]{\scriptsize High volatility}}%
\put(100,62){\makebox(61,5)[c]{\small (b) Absolute gamma error}}%
\put(100.000,22){\line(0,1){38}}%
\put(100.000,17){\makebox(0,3){\scriptsize 128}}%
{\color[HTML]{E5E5E5}\put(108.714,22){\line(0,1){38}}}%
{\color[HTML]{E5E5E5}\put(117.429,22){\line(0,1){38}}}%
\put(117.429,17){\makebox(0,3){\scriptsize 512}}%
{\color[HTML]{E5E5E5}\put(126.143,22){\line(0,1){38}}}%
{\color[HTML]{E5E5E5}\put(134.857,22){\line(0,1){38}}}%
\put(134.857,17){\makebox(0,3){\scriptsize 2048}}%
{\color[HTML]{E5E5E5}\put(143.571,22){\line(0,1){38}}}%
{\color[HTML]{E5E5E5}\put(152.286,22){\line(0,1){38}}}%
\put(152.286,17){\makebox(0,3){\scriptsize 8192}}%
\put(161.000,22){\line(0,1){38}}%
\put(161.000,17){\makebox(0,3){\scriptsize 16384}}%
{\color[HTML]{E5E5E5}\put(100.000,22.000){\line(1,0){61}}}%
\put(98.000,20.500){\makebox(0,3)[r]{\scriptsize $10^{-5}$}}%
{\color[HTML]{E5E5E5}\put(100.000,31.500){\line(1,0){61}}}%
\put(98.000,30.000){\makebox(0,3)[r]{\scriptsize $10^{-4}$}}%
{\color[HTML]{E5E5E5}\put(100.000,41.000){\line(1,0){61}}}%
\put(98.000,39.500){\makebox(0,3)[r]{\scriptsize $10^{-3}$}}%
{\color[HTML]{E5E5E5}\put(100.000,50.500){\line(1,0){61}}}%
\put(98.000,49.000){\makebox(0,3)[r]{\scriptsize $10^{-2}$}}%
{\color[HTML]{E5E5E5}\put(100.000,60.000){\line(1,0){61}}}%
\put(98.000,58.500){\makebox(0,3)[r]{\scriptsize $10^{-1}$}}%
\put(100,13){\makebox(61,3)[c]{\scriptsize Space intervals $M$}}%
{\color[HTML]{0072B2}\linethickness{0.5pt}\qbezier(100.000,47.306)(104.357,47.409)(108.714,47.513)}%
{\color[HTML]{0072B2}\linethickness{0.5pt}\qbezier(108.714,47.513)(113.071,46.206)(117.429,44.900)}%
{\color[HTML]{0072B2}\linethickness{0.5pt}\qbezier(117.429,44.900)(121.786,42.675)(126.143,40.449)}%
{\color[HTML]{0072B2}\linethickness{0.5pt}\qbezier(126.143,40.449)(130.500,40.749)(134.857,41.049)}%
{\color[HTML]{0072B2}\linethickness{0.5pt}\qbezier(134.857,41.049)(139.214,40.563)(143.571,40.077)}%
{\color[HTML]{0072B2}\linethickness{0.5pt}\qbezier(143.571,40.077)(147.929,38.847)(152.286,37.617)}%
{\color[HTML]{0072B2}\linethickness{0.5pt}\qbezier(152.286,37.617)(156.643,37.488)(161.000,37.359)}%
{\color[HTML]{0072B2}\put(100.000,47.306){\circle*{0.8}}}%
{\color[HTML]{0072B2}\put(108.714,47.513){\circle*{0.8}}}%
{\color[HTML]{0072B2}\put(117.429,44.900){\circle*{0.8}}}%
{\color[HTML]{0072B2}\put(126.143,40.449){\circle*{0.8}}}%
{\color[HTML]{0072B2}\put(134.857,41.049){\circle*{0.8}}}%
{\color[HTML]{0072B2}\put(143.571,40.077){\circle*{0.8}}}%
{\color[HTML]{0072B2}\put(152.286,37.617){\circle*{0.8}}}%
{\color[HTML]{0072B2}\put(161.000,37.359){\circle*{0.8}}}%
{\color[HTML]{0072B2}\put(100.000,8){\line(1,0){4}}}\put(105.000,6.5){\makebox(0,3)[l]{\scriptsize Total error}}%
{\color[HTML]{D55E00}\linethickness{0.5pt}\qbezier(100.000,50.957)(104.357,49.188)(108.714,47.418)}%
{\color[HTML]{D55E00}\linethickness{0.5pt}\qbezier(108.714,47.418)(113.071,45.986)(117.429,44.555)}%
{\color[HTML]{D55E00}\linethickness{0.5pt}\qbezier(117.429,44.555)(121.786,43.458)(126.143,42.362)}%
{\color[HTML]{D55E00}\linethickness{0.5pt}\qbezier(126.143,42.362)(130.500,41.536)(134.857,40.710)}%
{\color[HTML]{D55E00}\linethickness{0.5pt}\qbezier(134.857,40.710)(139.214,40.072)(143.571,39.434)}%
{\color[HTML]{D55E00}\linethickness{0.5pt}\qbezier(143.571,39.434)(147.929,38.915)(152.286,38.397)}%
{\color[HTML]{D55E00}\linethickness{0.5pt}\qbezier(152.286,38.397)(156.643,37.949)(161.000,37.501)}%
{\color[HTML]{D55E00}\put(100.000,50.957){\circle*{0.8}}}%
{\color[HTML]{D55E00}\put(108.714,47.418){\circle*{0.8}}}%
{\color[HTML]{D55E00}\put(117.429,44.555){\circle*{0.8}}}%
{\color[HTML]{D55E00}\put(126.143,42.362){\circle*{0.8}}}%
{\color[HTML]{D55E00}\put(134.857,40.710){\circle*{0.8}}}%
{\color[HTML]{D55E00}\put(143.571,39.434){\circle*{0.8}}}%
{\color[HTML]{D55E00}\put(152.286,38.397){\circle*{0.8}}}%
{\color[HTML]{D55E00}\put(161.000,37.501){\circle*{0.8}}}%
{\color[HTML]{D55E00}\put(132.000,8){\line(1,0){4}}}\put(137.000,6.5){\makebox(0,3)[l]{\scriptsize Bump bias}}%
{\color[HTML]{009E73}\linethickness{0.5pt}\qbezier(100.000,48.761)(104.357,40.326)(108.714,31.892)}%
{\color[HTML]{009E73}\linethickness{0.5pt}\qbezier(108.714,31.892)(113.071,39.741)(117.429,47.591)}%
{\color[HTML]{009E73}\linethickness{0.5pt}\qbezier(117.429,47.591)(121.786,42.931)(126.143,38.272)}%
{\color[HTML]{009E73}\linethickness{0.5pt}\qbezier(126.143,38.272)(130.500,34.421)(134.857,30.569)}%
{\color[HTML]{009E73}\linethickness{0.5pt}\qbezier(134.857,30.569)(139.214,31.330)(143.571,32.092)}%
{\color[HTML]{009E73}\linethickness{0.5pt}\qbezier(143.571,32.092)(147.929,31.614)(152.286,31.137)}%
{\color[HTML]{009E73}\linethickness{0.5pt}\qbezier(152.286,31.137)(156.643,27.335)(161.000,23.533)}%
{\color[HTML]{009E73}\put(100.000,48.761){\circle*{0.8}}}%
{\color[HTML]{009E73}\put(108.714,31.892){\circle*{0.8}}}%
{\color[HTML]{009E73}\put(117.429,47.591){\circle*{0.8}}}%
{\color[HTML]{009E73}\put(126.143,38.272){\circle*{0.8}}}%
{\color[HTML]{009E73}\put(134.857,30.569){\circle*{0.8}}}%
{\color[HTML]{009E73}\put(143.571,32.092){\circle*{0.8}}}%
{\color[HTML]{009E73}\put(152.286,31.137){\circle*{0.8}}}%
{\color[HTML]{009E73}\put(161.000,23.533){\circle*{0.8}}}%
{\color[HTML]{009E73}\put(100.000,4){\line(1,0){4}}}\put(105.000,2.5){\makebox(0,3)[l]{\scriptsize PDE difference}}%
\end{picture}}%
\caption{Synthetic convergence evidence. Panel (a) contrasts the lower and upper payoff kinks, a short-dated low-volatility claim and a high-volatility claim. Panel (b) separates the short-dated gamma total error into analytic bump bias and the PDE difference at the same bump. Absolute component errors need not add because their signed contributions can cancel. Both axes are logarithmic. All seven contracts and all five Greeks are retained in the comparison files.}
\label{fig:cgbm-convergence}
\end{figure}
Figure~\ref{fig:cgbm-convergence} exposes non-monotone refinement rather
than imposing a fitted rate. Kink/grid alignment and interpolation can
affect price error; differencing amplifies these errors in the Greeks.
For the short-dated gamma at the finest mesh, the analytic bump bias is
approximately $-4.28\times10^{-4}$, while the PDE difference at that bump
is approximately $1.45\times10^{-5}$. The larger gamma discrepancy is
therefore mainly a finite-bump effect. Finer price accuracy alone does not
remove it. Independent mesh refinements and all five Greek comparisons
remain available even where a selected plot would conceal adverse behaviour.

Every Monte Carlo method receives 262,144 terminal-payoff evaluations per
contract. Raw sampling uses that many independent normals. Antithetic
sampling uses 131,072 independent pairs. Controlled sampling uses 16,384
independent pilot pairs and 114,688 production pairs, charging both
terminal evaluations of each pair. A control coefficient is fitted
separately for each output and clipped to $[-10^6,10^6]$; the pilot is
independent of production. Shared production normals between methods are
declared, and no confidence claim is made for their paired differences.
The seed is 732451, with deterministic case offsets.
Terminal-evaluation budgets measure sampling work; arithmetic costs and
wall-clock time are not asserted to be identical between estimators.

The 49 raw outputs comprise the price, five Greeks and additional mixed
gamma in each of seven contracts. Each receives failure allowance
$0.05/49$. Their simultaneous finite-sample bounds apply to ideal
independent exact normal draws and exact arithmetic. Their floating-point
constants and pseudo-random normal generation have not been formally
certified. Normal sample-variance intervals are reported separately as
pointwise asymptotic diagnostics. Observed inclusion in the conservative
bounds is a sanity check, not a proof of coverage.

\begin{table}[tb]\centering\small
\begin{tabular}{lrrrr}\toprule
Contract & Analytic $\Gamma$ & Raw LR SE & Raw mixed SE & Variance ratio\\\midrule
Lower kink & $4.026\!\times\!10^{-3}$ & $7.689\!\times\!10^{-5}$ & $1.966\!\times\!10^{-5}$ & 15.293 \\
Upper kink & $-7.259\!\times\!10^{-3}$ & $6.210\!\times\!10^{-5}$ & $2.269\!\times\!10^{-5}$ & 7.493 \\
Inside spread & $-3.183\!\times\!10^{-3}$ & $7.019\!\times\!10^{-5}$ & $1.523\!\times\!10^{-5}$ & 21.240 \\
Deep cap & $-2.747\!\times\!10^{-3}$ & $4.014\!\times\!10^{-5}$ & $2.021\!\times\!10^{-5}$ & 3.945 \\
Low-volatility short dated & $0.296226$ & $3.320\!\times\!10^{-3}$ & $8.632\!\times\!10^{-4}$ & 14.792 \\
Negative rate/dividend & $2.713\!\times\!10^{-3}$ & $4.478\!\times\!10^{-5}$ & $1.272\!\times\!10^{-5}$ & 12.400 \\
High volatility & $-3.680\!\times\!10^{-4}$ & $4.251\!\times\!10^{-6}$ & $2.287\!\times\!10^{-6}$ & 3.455 \\
\bottomrule\end{tabular}
\caption{Raw gamma estimators at equal terminal-evaluation budgets. SE is the measured standard error, not a finite-sample confidence radius. The variance ratio is ordinary payoff-LR estimator variance divided by mixed-estimator variance. All antithetic and independently fitted control variants remain in the full comparison.}
\label{tab:cgbm-gamma-mc}\end{table}

The bounded mixed estimator reduces measured raw gamma variance in every
tested contract. This is an observed comparison, not uniform dominance
over model parameters. Controlled ordinary likelihood-ratio gamma remains
more efficient in several cases. Stock controls also worsen measured
delta and rho efficiency for the interior-spread contract at the matched
total budget. The analytic call spread has neither sampling error nor a
finite-domain or mesh error and remains the strongest baseline here.

\section{Interpretation and reproducibility}

This companion closes a specific preparation gap by connecting an
admissible financial strategy, a PDE, a nonsmooth numerical convergence
proof and independently checked sensitivities. The bounded mixed-gamma
derivation gives a useful estimator and a finite-sample uncertainty
statement for this claim. It is an explicit specialisation of sensitivity
principles, not a claimed new quantum theory or a completed priority claim.
The comparison retains negative gamma/vega, non-monotone errors and
adverse control rankings where they occur.

The study checks a stipulated one-asset model with known parameters. It
does not calibrate a financial model, validate market prices, solve a
general portfolio/control problem or establish a quantum advantage.
It also does not extend a one-asset analytic benchmark into a closed-form
solution for arbitrary baskets. Those boundaries are essential when
using the companion as evidence of financial mathematics and computation.

The engine is \texttt{capped\_gbm.py};
\texttt{run\_capped\_gbm\_study.py} reproduces the comparisons.
The tests include independent active-region Gaussian quadrature of the
mixed gamma, analytic sensitivity checks and the direct/sine-basis
backward-Euler comparison. The source ledger retains the final published
versions and precise sections consulted. Detailed arguments follow in
the appendices. I retain this document as a completed companion study;
the main quantum basket paper remains a working paper.

\appendix


\section{Financial and analytic derivations}
\label{app:cgbm-financial}

\subsection{Change of measure and the self-financing equation}

For a concrete construction, suppose under a physical measure $P$ that
$dS_t=(\mu-q)S_t\,dt+\sigma S_t\,dW_t^P$ with constant total-return drift
$\mu$. Set $\vartheta=(\mu-r)/\sigma$. Novikov's condition follows from
$\exp(\vartheta^2T/2)<\infty$ on every finite horizon. Thus
\begin{equation}
 \frac{dQ}{dP}\bigg|_{\mathcal F_T}
 =\exp\{-\vartheta W_T^P-\vartheta^2T/2\},\qquad
 W_t^Q=W_t^P+\vartheta t
 \label{eq:cgbm-rn}
\end{equation}
defines an equivalent probability measure and yields
\eqref{eq:cgbm-law}. The stochastic exponential is strictly positive
and has expectation one by the Gaussian moment-generating function.
For this constant coefficient, the Brownian conclusion can also be
checked directly. On any deterministic partition, the density is the
product of factors
$\exp\{-\vartheta\Delta W^P-\vartheta^2\Delta t/2\}$.
Multiplying the independent $N(0,\Delta t)$ increment densities by these
factors and completing the square changes each mean to
$-\vartheta\Delta t$, leaves its variance $\Delta t$ and preserves their
factorisation. Thus the increments of $W^Q=W^P+\vartheta t$ are independent
centred normals with the Brownian variances. The same calculation
conditional on the preceding Brownian filtration proves independence
from its past. Continuity is inherited from $W^P$. Hence $W^Q$ is a
Brownian motion under $Q$, an elementary verification of this established
constant-coefficient measure change.
Consequently this construction does not depend on an unverified
calibration of $\mu$.

Let $B_t=e^{rt}$, let $\delta_t$ be the share holding and let $\gamma_t$
be the bank-account holding. A self-financing wealth $P_t=\delta_tS_t+
\gamma_tB_t$ receives asset dividends, so
\begin{equation}
 dP_t=\delta_t\,dS_t+q\delta_tS_t\,dt+
 r(P_t-\delta_tS_t)\,dt
 =rP_t\,dt+\delta_t\sigma S_t\,dW_t^Q.
 \label{eq:cgbm-self-financing}
\end{equation}
For an admissible strategy, the share holding is predictable and
\begin{equation}
 \int_0^T e^{-2rt}\delta_t^2\sigma^2S_t^2\,dt<\infty
 \quad Q\text{-almost surely},\label{eq:cgbm-admissible-integrand}
\end{equation}
while discounted wealth is bounded below by a deterministic constant. For such a strategy, discounted
wealth is a local martingale bounded below and hence a supermartingale.
One may see the latter implication by localisation and conditional
Fatou's lemma after adding the fixed lower bound. An arbitrage with zero
initial wealth and non-negative, sometimes strictly positive terminal
wealth would contradict the supermartingale expectation inequality.
This proves no arbitrage within this explicit admissible strategy class;
it does not assert a general fundamental theorem for arbitrary markets.

\subsection{The analytic call and replication}

The event $X>A$ is $Z>-d_{2,A}$. Completing the square gives
\begin{equation}
 \mathbb E[e^{vZ}\mathbf 1_{\{Z>z\}}]
 =e^{v^2/2}\Phi(v-z).
 \label{eq:cgbm-gaussian-tilt}
\end{equation}
Substitution in $D\mathbb E[(X-A)^+]$ yields
\eqref{eq:cgbm-call}; payoff linearity yields the call spread. The normal
identity
\begin{equation}
 S e^{-q\tau}\varphi(d_{1,A})=A D\varphi(d_{2,A})
 \label{eq:cgbm-density-identity}
\end{equation}
follows directly by subtracting the two squared normal arguments. It
cancels the derivative-of-$\Phi$ terms when differentiating the call
formula. The normal arguments have derivatives
\begin{align}
 \partial_S d_{1,A}&=\partial_S d_{2,A}=\frac1{Sv},\notag\\
 \partial_\sigma d_{1,A}&=\sqrt{\tau}-\frac{d_{1,A}}{\sigma},
 &\partial_\sigma d_{2,A}&=-\frac{d_{1,A}}{\sigma},\notag\\
 \partial_r d_{1,A}&=\partial_r d_{2,A}=\frac{\sqrt{\tau}}{\sigma},
 \label{eq:cgbm-d-derivatives}\\
 \partial_\tau d_{1,A}&=
 \frac{r-q+\sigma^2/2}{\sigma\sqrt{\tau}}
 -\frac{d_{1,A}}{2\tau},
 &\partial_\tau d_{2,A}&=\partial_\tau d_{1,A}
 -\frac{\sigma}{2\sqrt{\tau}}.
 \notag
\end{align}
In the spot derivative of $C_A$, the two density terms cancel exactly,
leaving $e^{-q\tau}\Phi(d_{1,A})$; another spot derivative gives
$e^{-q\tau}\varphi(d_{1,A})/(Sv)$. In the volatility derivative, their
combined contribution is
$S e^{-q\tau}\varphi(d_{1,A})
(\partial_\sigma d_{1,A}-\partial_\sigma d_{2,A})
=S e^{-q\tau}\varphi(d_{1,A})\sqrt{\tau}$.
The rate-argument derivatives cancel, leaving
$\tau A D\Phi(d_{2,A})$ from differentiating the discount factor.
Finally,
\begin{equation}
 \partial_\tau C_A=-q S e^{-q\tau}\Phi(d_{1,A})
 +r A D\Phi(d_{2,A})
 +\frac{S e^{-q\tau}\sigma\varphi(d_{1,A})}{2\sqrt{\tau}}.
 \label{eq:cgbm-call-time}
\end{equation}
Subtracting the corresponding strike-$J$ expressions from those at
strike $K$ proves every formula in
\eqref{eq:cgbm-delta}--\eqref{eq:cgbm-theta}.
The lognormal convolution is smooth at positive maturity, even though
$F$ has two kinks. Dominated convergence proves continuity at expiry:
$X\to S$ almost surely, $F$ is continuous and $0\leq F\leq H$.

To verify the heat transformation explicitly, first write
$x=\log(S/K)$ and $V=e^{-r\tau}u(\tau,x)$. The chain rule gives
$SV_S=e^{-r\tau}u_x$ and
$S^2V_{SS}=e^{-r\tau}(u_{xx}-u_x)$. Equation
\eqref{eq:cgbm-pde} therefore reduces to
$u_\tau=a u_{xx}+b u_x$. Setting $u(\tau,x)=w(\tau,x+b\tau)$ gives
$w_\tau=a w_{yy}$. Conversely, Gaussian convolution proves
\eqref{eq:cgbm-heat}; reversing the calculation proves
\eqref{eq:cgbm-pde} without assuming it in the valuation argument.

With calendar time $t$ and expiry $T$, apply It\^o's formula to
$V(S_t,T-t)$ on $[0,T-\varepsilon]$. The pricing equation gives
\begin{equation}
 dV(S_t,T-t)=rV(S_t,T-t)\,dt+
 \sigma S_tV_S(S_t,T-t)\,dW_t^Q.
 \label{eq:cgbm-replication-ito}
\end{equation}
This agrees with \eqref{eq:cgbm-self-financing} for $\delta=V_S$ and
$\gamma=(V-SV_S)/B$. Furthermore
$0\leq V_S\leq e^{-q(T-t)}$ by \eqref{eq:cgbm-delta}, and
the Gaussian moment-generating function gives the explicit second moment
\begin{equation}
 \mathbb E_Q S_t^2=S_0^2
 \exp\{[2(r-q)+\sigma^2]t\}.
 \label{eq:cgbm-asset-second-moment}
\end{equation}
It is finite and bounded on every finite horizon. Hence
$\mathbb E_Q\int_0^T e^{-2rt}\sigma^2 S_t^2 V_S^2\,dt<\infty$.
The discounted stochastic gain is therefore a true square-integrable
martingale. Continuity at expiry gives terminal wealth $F(S_T)$ as
$\varepsilon\downarrow0$. Its discounted wealth is non-negative, so the
strategy is admissible. Any bounded quoted claim valuation different from
the replication cost permits buying the cheaper position, selling the dearer
one and investing the initial difference; their terminal claim cash
flows cancel. This gives a strictly positive terminal gain with zero
initial wealth and a bounded-below discounted wealth process, contradicting
no arbitrage. Any other bounded classical solution with this terminal
payoff has a discounted martingale representation with the same
conditional terminal expectation; equivalently, apply the stopped
It\^o argument and bounded convergence. Its price equals
\eqref{eq:cgbm-price}. This proves the stated uniqueness.

\section{Numerical stability and convergence proofs}
\label{app:cgbm-numerical}

\subsection{The matrix maximum principle}

Let $A_h$ be the $(M-1)\times(M-1)$ matrix with diagonal
$1+2\lambda$ and off-diagonals $-\lambda$. If $A_h z\geq0$ componentwise
and $z$ had a negative minimum at an interior row, that row would satisfy
$(A_hz)_j\leq z_j<0$. At an edge row, the missing neighbour is interpreted
as zero and gives an even stronger contradiction. Thus $z\geq0$.
The same argument applied to a null vector and its negative proves that
$A_h$ is nonsingular; consequently $A_h^{-1}\geq0$. The row sums of
$A_h$ are one in interior rows and $1+\lambda$ at the two edge rows.
Therefore $A_h\mathbf 1\geq\mathbf 1$, whence
\begin{equation}
 0\leq A_h^{-1}\mathbf 1\leq\mathbf 1,
 \qquad \|A_h^{-1}\|_\infty\leq1.
 \label{eq:cgbm-mmatrix}
\end{equation}
This proves contraction for two solutions with equal boundary values.
For positivity and the upper bound, the actual right-hand side is the
previous interior vector plus $\lambda H$ at the last row. If the previous
vector lies between $0$ and $H\mathbf 1$, this right-hand side lies between
$0$ and $A_h(H\mathbf 1)$. Applying the non-negative inverse proves the
invariant interval by induction. No condition relating $k$ and $h$ is used.

The discrete eigenvectors are $\sin(j\ell\pi/M)$ and substitution in the
centred second difference gives $-\nu_\ell$. The affine lift
$\ell(y)=H(y-L)/(U-L)$ has zero discrete second difference. Each lifted
mode is consequently multiplied by $(1+ak\nu_\ell)^{-1}$ per step,
which proves \eqref{eq:cgbm-dst} after $N$ steps. This identity is in exact
arithmetic; the direct-step comparison checks implementation error.

\subsection{Mollification without a false smooth-payoff assumption}

Let $\rho_\eta(z)=\eta^{-1}\rho(z/\eta)$ and
$f_\eta=f*\rho_\eta$. Global Lipschitz continuity gives
\begin{equation}
 \|f_\eta-f\|_\infty\leq L_fm_1\eta,
 \qquad
 \|f_\eta^{(4)}\|_\infty
 \leq L_f\|\rho_\eta'''\|_1=L_f c_3\eta^{-3}.
 \label{eq:cgbm-mollifier-bounds}
\end{equation}
The second estimate uses the weak first derivative $f'$, which exists
almost everywhere and is bounded by $L_f$. Moving the three remaining
derivatives to the smooth kernel proves the formula and its bound.
Since the radius is smaller than the distance of either kink to the
boundary, $f_\eta$ equals $0$ near $L$ and $H$ near $U$.

Subtract the affine lift from $f_\eta$ and shift the domain to
$[0,U-L]$. The lifted function is zero at both endpoints and affine in
their neighbourhoods. Its even derivatives, of order at least two,
vanish there. Odd reflection at both endpoints therefore produces a
smooth function of period $2(U-L)$: the value and every derivative agree
at each reflection. In particular, the fourth derivative of the extension
has maximum norm at most $L_f c_3\eta^{-3}$. Convolve this extension with
the whole-line Gaussian heat kernel and add the affine lift back. Oddness
keeps the lifted solution zero at both boundaries, so this is precisely
the finite-domain smooth solution $u_\eta$. Gaussian convolution contracts
the maximum norm and commutes with derivatives; hence
\begin{equation}
 \sup_{0\leq s\leq\tau}\|\partial_y^4u_\eta(s,\cdot)\|_\infty
 \leq L_f c_3\eta^{-3},\qquad
 \partial_s^2u_\eta=a^2\partial_y^4u_\eta.
 \label{eq:cgbm-smooth-regularity}
\end{equation}

At time $s_n=nk$, insert $u_\eta(s_n,y_j)$ into the numerical equation
and divide its residual by $k$. The integral remainder for backward
time differencing is at most
$k\sup|\partial_s^2u_\eta|/2$. The integral or Taylor remainder for the
centred second difference is at most
$h^2\sup|\partial_y^4u_\eta|/12$. Thus the residual $R_j^n$ obeys
\begin{equation}
 |R_j^n|\leq
 \left(\frac{a^2k}{2}+\frac{ah^2}{12}\right)L_f c_3\eta^{-3}.
 \label{eq:cgbm-residual}
\end{equation}
If the smooth grid starts from $f_\eta$, its error satisfies
$A_h e^n=e^{n-1}-kR^n$ with $e^0=0$. Equation
\eqref{eq:cgbm-mmatrix} and induction give
$\|e^N\|_\infty\leq Nk\sup_n\|R^n\|_\infty$.
Both the continuous maximum principle and the discrete contraction bound
the change from $f_\eta$ back to $f$ by $L_fm_1\eta$. Adding these two
changes proves \eqref{eq:cgbm-be-error}. Taking
$\eta$ proportional to $(h^2+k)^{1/4}$, eventually below the fixed boundary
margin, proves convergence and the stated upper rate. This is a proof
for the actual kinked initial condition, not for a replacement smooth
contract.

\subsection{Domain truncation and interpolation}

Let $Y_s=x+\sigma B_s$, and let $\zeta$ be its first exit from $[L,U]$.
The stopped Brownian representation of the finite-domain solution is
\begin{equation}
 w^{\mathrm{box}}(\tau,x)=
 \mathbb E_x[f(Y_\tau)\mathbf 1_{\{\zeta>\tau\}}]
 +\mathbb E_x[g(Y_\zeta)\mathbf 1_{\{\zeta\leq\tau\}}],
 \qquad g(L)=0,\quad g(U)=H.
 \label{eq:cgbm-stopped-heat}
\end{equation}
It follows either from the stopped It\^o formula for positive-time smooth
solutions and bounded convergence, or from the killed heat semigroup.
The full solution is $\mathbb E_x f(Y_\tau)$. On exit, the terminal payoff
and boundary value are both in $[0,H]$, so their difference is bounded
in absolute value by $H$. The reflection principle bounds each one-sided
exit event by twice its corresponding normal tail; a union bound proves
\eqref{eq:cgbm-domain-beta}. The upper and lower exit contributions can
have opposite signs, so this is an absolute error bound rather than an
ordering of the two prices.

At positive maturity, differentiate the whole-line convolution once using
the bounded weak derivative $f'$, then once using the Gaussian density.
For a Gaussian with standard deviation $v$, the $L^1$ norm of its density
derivative is $\sqrt{2/\pi}/v$. Therefore
\begin{equation}
 \|w_{yy}(\tau,\cdot)\|_\infty
 \leq \frac{L_f}{\sigma\sqrt{\tau}}\sqrt{\frac{2}{\pi}}.
 \label{eq:cgbm-heat-curvature}
\end{equation}
Linear interpolation on an interval of width $h$ has error at most
$h^2\sup|w_{yy}|/8$. Interpolation weights are non-negative and sum to
one. Apply the domain bound at the two bracketing nodes and the nodal
discretisation bound to these same weights; add the whole-line
interpolation error and discount. This proves
\eqref{eq:cgbm-total-error}, without requiring a new boundary regularity
bound for $w^{\mathrm{box}}$.

\subsection{Uniform consistency for Greeks}

For a central first difference, price error contributes at most
$(E+E)/(2d)=E/d$. Taylor's theorem with an integral remainder gives
bias at most $M_3d^2/6$. For the central second difference, price error
contributes at most $(E+2E+E)/d^2=4E/d^2$, and the fourth-order remainder
is at most $M_4d^2/12$. These prove
\eqref{eq:cgbm-first-bump} and \eqref{eq:cgbm-second-bump}.

On a compact parameter set with positive lower bounds for $S$, $\sigma$
and $\tau$, the analytic formula is smooth and its third and fourth
parameter derivatives are bounded. Its transformed queries also lie in
a compact interval. Let $L\to-\infty$ and $U\to\infty$, keep
$h\to0$ and $k\to0$, and use a radius proportional to
$(h^2+k)^{1/4}$. The Gaussian exit tails vanish uniformly because
$\sigma\sqrt{\tau}$ has a finite upper bound, and the other terms in
\eqref{eq:cgbm-total-error} vanish uniformly. The discount factors are
bounded even when $r$ is negative. Thus $E_n\to0$ uniformly. Choosing
the bump sequence as stated in Proposition~\ref{prop:cgbm-bumps} proves
Greek convergence. A fixed domain or a fixed price error does not meet
this argument merely because its bumps are made smaller.

\section{Monte Carlo derivations and uncertainty proofs}
\label{app:cgbm-monte-carlo}

\subsection{Pathwise first derivatives}

The capped payoff is Lipschitz in its asset argument, with derivative
$\mathbf 1_{\{K<x<J\}}$ except at two points. The continuously distributed
$X$ hits those points with probability zero. Its parameter derivatives are
\begin{equation}
 X_S=X/S,\qquad X_\sigma=X(\sqrt{\tau}Z-\sigma\tau),\qquad
 X_r=\tau X,\qquad
 X_\tau=X\left(b+\frac{\sigma Z}{2\sqrt{\tau}}\right).
 \label{eq:cgbm-path-derivatives}
\end{equation}
On any compact parameter neighbourhood with $S,\sigma,\tau>0$, difference
quotients of the Lipschitz payoff are dominated by a constant times
$(1+|Z|)\exp(C|Z|)$, which is integrable against a Gaussian. This follows
from the mean-value theorem for the smooth terminal map $X$ and the
Lipschitz bound for $F$. Dominated convergence therefore justifies
interchanging differentiation and expectation. Differentiating the
discount factor as well gives
\eqref{eq:cgbm-pw-estimators} and \eqref{eq:cgbm-theta-estimator}.
Neither payoff kink contributes a first-order term because the payoff
itself is continuous there.

On the active region, $X\leq J$ and $Z\in(z_-,z_+)$, giving the delta,
vega and theta ranges in \eqref{eq:cgbm-mc-ranges}. For rho the expression
$IX-F(X)$ is exactly $0$, $K$ or $-H$ in the below-strike, active and
capped regions, respectively. This proves its stated tighter range.

\subsection{The likelihood-ratio gamma}

For $x>0$, the lognormal density is
\begin{equation}
 p_S(x)=\frac{1}{xv\sqrt{2\pi}}
 \exp\!\left[-\frac{1}{2}
 \left(\frac{\log(x/S)-b\tau}{v}\right)^2\right].
 \label{eq:cgbm-density}
\end{equation}
Writing $z=(\log(x/S)-b\tau)/v$, direct differentiation gives
\begin{equation}
 \frac{\partial_Sp_S(x)}{p_S(x)}=\frac{z}{Sv},\qquad
 \frac{\partial_S^2p_S(x)}{p_S(x)}
 =\frac{1}{S^2}\left(\frac{z^2-1}{v^2}-\frac{z}{v}\right).
 \label{eq:cgbm-lr-weight}
\end{equation}
On a compact spot neighbourhood away from zero, the density and its first
two spot derivatives admit a common envelope
\begin{equation}
 C(1+|\log x|^2)x^{-1}
 \exp\{-|\log x|^2/(4v^2)\},\qquad x>0,
 \label{eq:cgbm-lr-envelope}
\end{equation}
for a finite constant $C$. Indeed $\log S+b\tau$ is bounded on the
neighbourhood and $(u-m)^2\geq u^2/2-m^2$; the derivative weights are
polynomials of degree at most two in $u=\log x$, with bounded inverse
spot factors. The envelope is integrable by the substitution $u=\log x$,
and remains integrable after multiplication by $F\leq H$.
Two applications of dominated convergence
to $D\int F(x)p_S(x)\,dx$ prove unbiasedness of
\eqref{eq:cgbm-lr-estimator}. Using $F\leq H$ gives
\begin{align}
 \mathbb E[Y_\Gamma^2]
 &\leq \frac{D^2H^2}{S^4}
 \mathbb E\left[
 \left(\frac{Z^2-1}{v^2}-\frac{Z}{v}\right)^2\right]\notag\\
 &=\frac{D^2H^2}{S^4}
 \left(\frac{2}{v^4}+\frac{1}{v^2}\right),
 \label{eq:cgbm-lr-moment}
\end{align}
because $\mathbb E(Z^2-1)^2=2$, $\mathbb E Z^2=1$ and the odd mixed
moment is zero. Variance is no larger than the second moment, proving
\eqref{eq:cgbm-gamma-variance}. Gamma's likelihood-ratio estimator is
unbounded; a bounded-payoff concentration inequality cannot be applied
to it without this separate argument.

\subsection{A bounded mixed gamma rival}

The pathwise delta has the fixed-asset representation
\begin{equation}
 \Delta=\frac{D}{S}\int_K^J x p_S(x)\,dx.
 \label{eq:cgbm-delta-density}
\end{equation}
The integration limits are fixed when spot varies. On this bounded
interval away from zero, $x\partial_S p_S(x)$ has a common integrable
envelope on every compact spot neighbourhood. Differentiation is
therefore legitimate even though the payoff derivative is an indicator.
Using \eqref{eq:cgbm-lr-weight},
\begin{align}
 \partial_S\Delta
 &=-\frac{D}{S^2}\int_K^J x p_S(x)\,dx
 +\frac{D}{S}\int_K^J x\partial_Sp_S(x)\,dx\notag\\
 &=\frac{D}{S^2}\int_K^J x p_S(x)
 \left(\frac{z}{v}-1\right)\,dx.
 \label{eq:cgbm-mixed-gamma-derivation}
\end{align}
This proves unbiasedness of \eqref{eq:cgbm-mixed-gamma}. It includes the
boundary effect through the density derivative. On the active region,
$X<J$ and $|Z|\leq m$, proving \eqref{eq:cgbm-mixed-gamma-bound} and its
finite second moment. As an independent analytic check, Gaussian tilting
reduces its expectation to
$e^{-q\tau}\{\varphi(d_{1,K})-\varphi(d_{1,J})\}/(Sv)$,
in agreement with \eqref{eq:cgbm-delta-gamma}. I combine the established
pathwise and likelihood-ratio constructions of
\citet{broadieglasserman1996}, deriving the capped-payoff expression
explicitly here; I make no priority claim. The ordinary and mixed gamma
estimators are compared using identical financial parameters and matched
terminal evaluation budgets; neither estimator is discarded after
seeing its variance.

\subsection{Confidence guarantees and stronger classical sampling}

For completeness, the bounded-variable concentration argument of
\citet{hoeffding1963} can be checked directly. If $Y$ lies in an interval
of width $R>0$, define
$A(s)=\log\mathbb E\exp\{s(Y-\mathbb E Y)\}$ for real $s$.
Boundedness justifies differentiating twice. The second derivative is
the variance of $Y$ under the exponentially tilted probability law.
That law has the same interval support; centring at its midpoint gives
$A''(s)\leq R^2/4$. Since $A(0)=A'(0)=0$, integration twice yields
\begin{equation}
 A(s)\leq R^2s^2/8,\qquad s\in\mathbb R.
 \label{eq:cgbm-hoeffding-mgf}
\end{equation}
For independent observations and $s>0$, exponential Markov's inequality
and the product moment-generating function therefore give
\begin{equation}
 \mathbb P(\overline Y-\mathbb E Y\geq t)
 \leq\exp\{n(R^2s^2/8-st)\}.
 \label{eq:cgbm-hoeffding-chernoff}
\end{equation}
The minimiser is $s=4t/R^2$, giving the upper-tail bound
$\exp(-2nt^2/R^2)$. Apply the same argument to $-Y$ and sum the two
tail bounds. Solving
$2\exp(-2nt^2/R^2)=\alpha$ proves
\eqref{eq:cgbm-hoeffding}. If $R=0$, the observation is constant and
its confidence radius is zero. Independence gives
$\operatorname{Var}(\overline Y_\Gamma)\leq B_\Gamma^2/n$;
Chebyshev's inequality gives the stated gamma radius. A union bound
requires no independence between different outputs calculated on the
same paths. Their allocated failure probabilities must sum to the
intended familywise level.

For the stock control $C=D X-S e^{-q\tau}$, the Gaussian exponential
moment gives $\mathbb E C=0$ and a finite second moment. For any estimator
$Y$, a deterministic coefficient $\beta$ yields
\begin{equation}
 \operatorname{Var}(Y-\beta C)
 =\operatorname{Var}(Y)-2\beta\operatorname{Cov}(Y,C)
 +\beta^2\operatorname{Var}(C).
 \label{eq:cgbm-control-variance}
\end{equation}
When $\operatorname{Var}(C)>0$, the population minimiser is
$\beta_* =\operatorname{Cov}(Y,C)/\operatorname{Var}(C)$ and the
minimum variance is
$\operatorname{Var}(Y)-\operatorname{Cov}(Y,C)^2/
\operatorname{Var}(C)$. An estimated pilot coefficient need not attain
that minimum. Independence of pilot and main sample preserves conditional
unbiasedness; a zero pilot denominator is handled by selecting zero
control coefficient. Antithetic pairs replace $Y(Z)$ and $C(Z)$ with
their averages at $Z$ and $-Z$ and preserve their expectations, but variance
reduction for every Greek is not guaranteed merely by symmetry. The
comparison therefore reports measured variance and complete sampling work.
For fixed finite pilot coefficients, the main-pair estimators have finite
variance, so their studentised normal intervals have asymptotic validity
under non-zero variance. Finite-sample exact coverage is not claimed.

\begin{thebibliography}{99}
\bibitem[Black and Scholes(1973)]{blackscholes1973}
Black, F. and Scholes, M. (1973). The pricing of options and corporate liabilities. \emph{Journal of Political Economy}, 81(3), 637--654. \url{https://doi.org/10.1086/260062}. Final journal scan consulted, printed pages 637--644.

\bibitem[Merton(1973)]{merton1973}
Merton, R. C. (1973). Theory of rational option pricing. \emph{The Bell Journal of Economics and Management Science}, 4(1), 141--183. \url{https://doi.org/10.2307/3003143}. Final journal PDF consulted, selected material in Sections 6--7, printed pages 162--166 and 168--171, particularly equation (44), page 170, and footnote 62, page 171.

\bibitem[Broadie and Glasserman(1996)]{broadieglasserman1996}
Broadie, M. and Glasserman, P. (1996). Estimating security price derivatives using simulation. \emph{Management Science}, 42(2), 269--285. \url{https://doi.org/10.1287/mnsc.42.2.269}. Final journal scan consulted, Sections 2--3, printed pages 270--274, and Appendices A--C, printed pages 280--284.

\bibitem[Giles and Carter(2006)]{gilescarter2006}
Giles, M. B. and Carter, R. (2006). Convergence analysis of Crank--Nicolson and Rannacher time-marching. \emph{Journal of Computational Finance}, 9(4), 89--112. \url{https://doi.org/10.21314/JCF.2006.152}. Final journal PDF consulted, Sections 1--3, printed pages 89--103, and Section 5, page 111.

\bibitem[Hoeffding(1963)]{hoeffding1963} Hoeffding, W. (1963). Probability inequalities for sums of bounded random variables. \emph{Journal of the American Statistical Association}, 58(301), 13--30. \url{https://doi.org/10.1080/01621459.1963.10500830}. Earlier final-journal-scan consultation retained in the project source ledger, printed pages 13--16, particularly Theorem 2, equation (2.6).

\end{thebibliography}
\end{document}
